\section{Results: protease (diversity-reward ablation)}
\label{sec:res:protease}

\plan{\textbf{Claim.} A batch-consensus mutation-diversity reward accesses
rarer mutations and improves catalytic activity over an otherwise-matched
library that lacks the diversity term. \textbf{Risk --- highest in the paper.}
It depends on wet-lab activity data for both arms. The load-bearing finding is
``rare mutations \emph{improve} activity,'' not merely ``we sampled more of
sequence space''; without that activity uplift this section reduces to a
diversity-for-diversity's-sake observation. Status: \statusblocked.}

\subsection*{Block 1 --- Setup}
\begin{planlist}
  \item De novo protease backbone.
  \item Two arms, matched library budget: (i) RL fine-tuned MPNN, standard
        reward set; (ii) same as (i) plus the batch-consensus
        mutation-diversity reward
        (\texttt{cleo.design.utils.mutation\_diversity}).
\end{planlist}
\todo{protease results --- setup, both arms}

\subsection*{Block 2 --- Sequence-space coverage}
\begin{planlist}
  \item Mutation-frequency spectra: arm (ii) should populate rarer mutations
        more heavily than arm (i).
  \item Position-level entropy / consensus-divergence statistics.
\end{planlist}
\todo{protease results --- sequence-space coverage figure}

\subsection*{Block 3 --- Activity comparison}
\begin{planlist}
  \item Activity distribution: arm (i) vs.\ arm (ii).
  \item Best-variant comparison.
  \item Are the rare-mutation-bearing variants the ones with highest activity?
        This is the punchline --- call it out explicitly.
\end{planlist}
\todo{protease results --- activity comparison; rare-mutation/activity link}

\subsection*{Block 4 --- Mechanism for the diversity payoff}
\begin{planlist}
  \item Which rare mutations contributed; structural rationalization if
        possible.
\end{planlist}
\todo{protease results --- mechanistic rationalization}
